%============================================================
\chapter{検証障害の分類：六つの状態}
\label{part:obstacles}
%============================================================

\section{分類の必要}

「検証できない」には性質の異なる状態が含まれる。
対処法が異なるため、分けて扱う必要がある。この分類は本稿の整理である。

第\ref{ch:design}章の五段階のうち、段階3（取得可能性）で詰まったとき、
何が詰まっているのかを言い分けられなければ、
労力を投じるべきか断念すべきかを決められない。

\section{六分類}
\label{sec:obstacle-taxonomy}

\begin{center}
\begin{tabularx}{\textwidth}{@{}l>{\raggedright\arraybackslash}p{0.36\textwidth}X@{}}
\toprule
種別 & 状態 & 正しい対処 \\
\midrule
（i）測定の不在 & 世界の誰もその量を測っていない & 自ら原調査を行う。二次データでは永久に埋まらない \\
（ii）仮説の不良設定 & 理論が符号を一意に決めない & 理論を精密化する。データを増やしても解決しない \\
（iii）識別不能 & データはあるが交絡を切れない & 設計の工夫。外生変動か対照群を探す \\
（iv）アクセス制約 & 存在するが取得できない & 経路の開拓、または法的・技術的に断念 \\
（v）集計の選択 & 調査はされるが公表される集計表に含まれない & 集計要望、匿名データの利用申請 \\
（vi）調整の不作動 & データは十分にあるが、説明変数が理論の想定する範囲で変動しない & \textbf{どの経路でも解決しない} \\
\bottomrule
\end{tabularx}
\captionof{table}{検証障害の六分類。状態と正しい対処。}
\label{tab:obstacle-taxonomy}
\label{sec:sixth-obstacle}
\end{center}

\textbf{労力で動くものと動かないものを分けるのが、この分類の役割である。}

（iii）（iv）（v）は労力や工夫で動きうる。
（iii）は設計、（iv）は経路、（v）は公表の範囲がそれぞれ変われば解決する。

（i）（ii）（vi）は動かない。
（i）は測定そのものを起こさない限り埋まらず、
（ii）はデータを増やしても符号が決まらず、
（vi）は説明変数が動いていない以上、いくら標本を増やしても検定が成立しない。

\begin{remark}[（vi）は否定的な結果ではない]
\label{rem:vi-is-observation}
（vi）は理論の欠陥でも測定の欠陥でもない。
\textbf{理論が想定する調整機構が現実に作動していないという事実}である。

理論の予測が現れないことは、それ自体が観測である。
第\ref{part:results}部では「手数料率が成果のリスクに応じて調整されていない」
という記述がこれにあたる（注\ref{rem:no-adjustment}）。

第\ref{sec:layers}節のレイヤー構造でいえば、
価格が慣行で決まることは制度層に属する現象である。
最適化を前提とする理論が、慣行の前で作動しないという構図になる。
\end{remark}

\begin{remark}[集計の不能は集計の選択と異なる]
\label{rem:agg-impossible}
（v）は集計されていないだけであり、標本には含まれている。
これに対し、観測の単位が主体と対応しないために
主体単位の量を復元できない場合がある。
第\ref{ch:protocol}章の識別子がこれにあたる（注\ref{rem:additivity}）。

こちらは公表の範囲が変わっても解決しない。
（v）ではなく（i）に近いが、測定されていないのではなく
\textbf{測定の単位が問いの単位と違う}という点で別である。
\end{remark}

\section{分類は着手前の判定にも使える}

この分類は、検証が完結しなかった原因の事後の整理であると同時に、
\textbf{調査に着手する前の判定にも使える}。

第\ref{ch:design}章の段階3において、
当該の量が必要な粒度で取得できるかを判定する際の基準がこれにあたる。
第\ref{ch:platformfee}章の調査可能性の判定は、この使い方の実例である。

%============================================================
